% Dokumentenkopf ---------------------------------------------------------------
%   Diese Vorlage basiert auf "scrreprt" aus dem koma-script.
% ------------------------------------------------------------------------------
\documentclass[
    12pt, % Schriftgröße
    DIV10,
    ngerman, % für Umlaute, Silbentrennung etc.
    a4paper, % Papierformat
    oneside, % einseitiges Dokument
    notitlepage, % es wird eine Titelseite verwendet, aber mit Abstract
    parskip=half, % Abstand zwischen Absätzen (halbe Zeile)
    headings=normal, % Größe der Überschriften verkleinern
    listof=totoc, % Verzeichnisse im Inhaltsverzeichnis aufführen
    bibliography=totoc, % Literaturverzeichnis im Inhaltsverzeichnis aufführen
    index=totoc, % Index im Inhaltsverzeichnis aufführen
    captions=tableheading, % Beschriftung von Tabellen unterhalb ausgeben
    final % Status des Dokuments (final/draft)
]{scrreprt}

% Meta-Informationen -----------------------------------------------------------
%   Definition von globalen Parametern, die im gesamten Dokument verwendet
%   werden k�nnen (z.B auf dem Deckblatt etc.).
%
%   ACHTUNG: Wenn die Texte Umlaute oder ein Esszet enthalten, muss der folgende
%            Befehl bereits an dieser Stelle aktiviert werden:
%            \usepackage[latin1]{inputenc}
% ------------------------------------------------------------------------------
\newcommand{\titel}{Clojure}
\newcommand{\untertitel}{und hier kommt der Untertitel}
\newcommand{\art}{Seminararbeit}
\newcommand{\fachgebiet}{Informatik}
\newcommand{\autor}{Benjamin Britsch, Alain M. Lafon}
\newcommand{\studienbereich}{Informatik}
\newcommand{\matrikelnr}{373679, 372991}
\newcommand{\erstgutachter}{Prof. Dr. Stefan Knauth}
\newcommand{\jahr}{2011}
\newcommand{\ort}{Stuttgart, Z�rich}
\newcommand{\logo}{Bilder/Hft_Logo_Klein.png}

% Anpassung des Seitenlayouts --------------------------------------------------
%   siehe Seitenstil.tex
% ------------------------------------------------------------------------------
\usepackage[
    automark, % Kapitelangaben in Kopfzeile automatisch erstellen
    headsepline, % Trennlinie unter Kopfzeile
    ilines % Trennlinie linksb�ndig ausrichten
]{scrpage2}

% Anpassung an Landessprache ---------------------------------------------------
\usepackage[ngerman]{babel}

% Umlaute ----------------------------------------------------------------------
%   Umlaute/Sonderzeichen wie ���� direkt im Quelltext verwenden (CodePage).
%   Erlaubt automatische Trennung von Worten mit Umlauten.
% ------------------------------------------------------------------------------
\usepackage[latin1]{inputenc}
\usepackage[T1]{fontenc}
\usepackage{textcomp} % Euro-Zeichen etc.

% Schrift ----------------------------------------------------------------------
\usepackage{lmodern} % bessere Fonts
\usepackage{relsize} % Schriftgr��e relativ festlegen

% Grafiken ---------------------------------------------------------------------
% Einbinden von JPG-Grafiken erm�glichen
\usepackage[dvips,final]{graphicx}
% hier liegen die Bilder des Dokuments
\graphicspath{{Bilder/}}

% Befehle aus AMSTeX f�r mathematische Symbole z.B. \boldsymbol \mathbb --------
\usepackage{amsmath,amsfonts}

% f�r Index-Ausgabe mit \printindex --------------------------------------------
\usepackage{makeidx}

% Einfache Definition der Zeilenabst�nde und Seitenr�nder etc. -----------------
\usepackage{setspace}
\usepackage{geometry}

% Symbolverzeichnis ------------------------------------------------------------
%   Symbolverzeichnisse bequem erstellen. Beruht auf MakeIndex:
%     makeindex.exe %Name%.nlo -s nomencl.ist -o %Name%.nls
%   erzeugt dann das Verzeichnis. Dieser Befehl kann z.B. im TeXnicCenter
%   als Postprozessor eingetragen werden, damit er nicht st�ndig manuell
%   ausgef�hrt werden muss.
%   Die Definitionen sind ausgegliedert in die Datei "Glossar.tex".
% ------------------------------------------------------------------------------
\usepackage[intoc]{nomencl}
\let\abbrev\nomenclature
\renewcommand{\nomname}{Abk�rzungsverzeichnis}
\setlength{\nomlabelwidth}{.25\hsize}
\renewcommand{\nomlabel}[1]{#1 \dotfill}
\setlength{\nomitemsep}{-\parsep}

% zum Umflie�en von Bildern ----------------------------------------------------
\usepackage{floatflt}


% zum Einbinden von Programmcode -----------------------------------------------
\usepackage{listings}
\usepackage{xcolor} 
\definecolor{hellgelb}{rgb}{1,1,0.9}
\definecolor{colKeys}{rgb}{0,0,1}
\definecolor{colIdentifier}{rgb}{0,0,0}
\definecolor{colComments}{rgb}{1,0,0}
\definecolor{colString}{rgb}{0,0.5,0}
\lstset{
    float=hbp,
    basicstyle=\ttfamily\color{black}\small\smaller,
    identifierstyle=\color{colIdentifier},
    keywordstyle=\color{colKeys},
    stringstyle=\color{colString},
    commentstyle=\color{colComments},
    columns=flexible,
    tabsize=2,
    frame=single,
    extendedchars=true,
    showspaces=false,
    showstringspaces=false,
    numbers=left,
    numberstyle=\tiny,
    breaklines=true,
    backgroundcolor=\color{hellgelb},
    breakautoindent=true
}

% URL verlinken, lange URLs umbrechen etc. -------------------------------------
\usepackage{url}

% wichtig f�r korrekte Zitierweise ---------------------------------------------
\usepackage[square]{natbib}

% PDF-Optionen -----------------------------------------------------------------
\usepackage[
    bookmarks,
    bookmarksopen=true,
    colorlinks=true,
% diese Farbdefinitionen zeichnen Links im PDF farblich aus
    linkcolor=red, % einfache interne Verkn�pfungen
    anchorcolor=black,% Ankertext
    citecolor=blue, % Verweise auf Literaturverzeichniseintr�ge im Text
    filecolor=magenta, % Verkn�pfungen, die lokale Dateien �ffnen
    menucolor=red, % Acrobat-Men�punkte
    urlcolor=cyan, 
% diese Farbdefinitionen sollten f�r den Druck verwendet werden (alles schwarz)
    %linkcolor=black, % einfache interne Verkn�pfungen
    %anchorcolor=black, % Ankertext
    %citecolor=black, % Verweise auf Literaturverzeichniseintr�ge im Text
    %filecolor=black, % Verkn�pfungen, die lokale Dateien �ffnen
    %menucolor=black, % Acrobat-Men�punkte
    %urlcolor=black, 
    backref,
    plainpages=false, % zur korrekten Erstellung der Bookmarks
    pdfpagelabels, % zur korrekten Erstellung der Bookmarks
    hypertexnames=false, % zur korrekten Erstellung der Bookmarks
    linktocpage % Seitenzahlen anstatt Text im Inhaltsverzeichnis verlinken
]{hyperref}
% Befehle, die Umlaute ausgeben, f�hren zu Fehlern, wenn sie hyperref als Optionen �bergeben werden
\hypersetup{
    pdftitle={\titel \untertitel},
    pdfauthor={\autor},
    pdfcreator={\autor},
    pdfsubject={\titel \untertitel},
    pdfkeywords={\titel \untertitel},
}

% fortlaufendes Durchnummerieren der Fu�noten ----------------------------------
\usepackage{chngcntr}

% f�r lange Tabellen -----------------------------------------------------------
\usepackage{longtable}
\usepackage{array}
\usepackage{ragged2e}
\usepackage{lscape}

% Spaltendefinition rechtsb�ndig mit definierter Breite ------------------------
\newcolumntype{w}[1]{>{\raggedleft\hspace{0pt}}p{#1}}

% Formatierung von Listen �ndern -----------------------------------------------
\usepackage{paralist}

% bei der Definition eigener Befehle ben�tigt
\usepackage{ifthen}

% definiert u.a. die Befehle \todo und \listoftodos
\usepackage{todonotes}

% sorgt daf�r, dass Leerzeichen hinter parameterlosen Makros nicht als Makroendezeichen interpretiert werden
\usepackage{xspace}

% CTAN http://tug.ctan.org/cgi-bin/ctanPackageInformation.py?id=caption
% Implementiert beispielsweise Umbr�che in Tabellen-Beschriftungen
\usepackage[format=plain,justification=centerlast ]{caption}


% Index und Abkürzungsverzeichnisses
\makeindex
\makenomenclature


% Kopf- und Fußzeilen, Seitenränder, etc
% Zeilenabstand 1,5 Zeilen -----------------------------------------------------
\onehalfspacing

% Seitenr�nder -----------------------------------------------------------------
\setlength{\topskip}{\ht\strutbox} % behebt Warnung von geometry
\geometry{paper=a4paper,left=35mm,right=35mm,top=30mm}

% Kopf- und Fu�zeilen ----------------------------------------------------------
\pagestyle{scrheadings}
% Kopf- und Fu�zeile auch auf Kapitelanfangsseiten
\renewcommand*{\chapterpagestyle}{scrheadings} 
% Schriftform der Kopfzeile
\renewcommand{\headfont}{\normalfont}

% Kopfzeile
\ihead{\large{\textsc{\titel}}\\ [2ex] \textit{\headmark}}
\chead{}
\ohead{\includegraphics[scale=0.4]{\logo}}
\setlength{\headheight}{21mm} % H�he der Kopfzeile
% Kopfzeile �ber den Text hinaus verbreitern
\setheadwidth[0pt]{textwithmarginpar} 
\setheadsepline[text]{0.4pt} % Trennlinie unter Kopfzeile

% Fu�zeile
%\ifoot{\copyright\ \autor}
\cfoot{}
\ofoot{\pagemark}

% sonstige typographische Einstellungen ----------------------------------------

% erzeugt ein wenig mehr Platz hinter einem Punkt
\frenchspacing 

% Schusterjungen und Hurenkinder vermeiden
\clubpenalty = 10000
\widowpenalty = 10000 
\displaywidowpenalty = 10000

% Quellcode-Ausgabe formatieren
\lstset{numbers=left, numberstyle=\tiny, numbersep=5pt, breaklines=true}
\lstset{emph={square}, emphstyle=\color{red}, emph={[2]root,base}, emphstyle={[2]\color{blue}}}

% Fu�noten fortlaufend durchnummerieren
\counterwithout{footnote}{chapter}


% eigene Definitionen für Silbentrennung
% Trennvorschl�ge im Text werden mit \" angegeben
% untrennbare W�rter und Ausnahmen von der normalen Trennung k�nnen in dieser
% Datei mittels \hyphenation definiert werden


% eigene LaTeX-Befehle
% Eigene Befehle und typographische Auszeichnungen f�r diese

% Abk�rzungen mit korrektem Leerraum 
\newcommand{\ua}{\mbox{u.\,a.\ }}
\newcommand{\zB}{\mbox{z.\,B.\ }}
\newcommand{\dahe}{\mbox{d.\,h.\ }}
\newcommand{\Vgl}{Vgl.\ }
\newcommand{\bzw}{bzw.\ }
\newcommand{\evtl}{evtl.\ }

\newcommand{\bs}{$\backslash$}

% zum Ausgeben von Autoren
\newcommand{\AutorName}[1]{\textsc{#1}}
\newcommand{\Autor}[1]{\AutorName{\citeauthor{#1}}}

% verschiedene Befehle um W�rter semantisch auszuzeichnen 
\newcommand{\NeuerBegriff}[1]{\textbf{#1}}
\newcommand{\Fachbegriff}[1]{\emph{#1}}

\newcommand{\Eingabe}[1]{\texttt{#1}}
\newcommand{\Code}[1]{\texttt{#1}}
\newcommand{\Datei}[1]{\texttt{#1}}

\newcommand{\Datentyp}[1]{\textsf{#1}}


\begin{document}

% erlaube regulaere Anfuehrungszeichen
\shorthandoff{"}

% subsubsection nummerieren
\setcounter{secnumdepth}{3}
\setcounter{tocdepth}{3}

% Deckblatt und Abstract ohne Seitenzahl
\ofoot{}
\thispagestyle{plain}
\begin{titlepage}

\begin{center}

\LARGE{{\titel}}\\[1.5ex]
%\huge{\textbf{\titel}}\\[1.5ex]
%\LARGE{\textbf{\untertitel}}\\[6ex]
%\LARGE{\textbf{\art}}\\[1.5ex]
%\includegraphics[scale=0.7]{Hft_Logo.png}\\%\\[6ex]
%\includegraphics[scale=0.5]{Swiss_Post_Logo_inverted.png}%\\[6ex]

\end{center}
\footnotesize
\begin{tabular}{w{5cm}p{6cm}}\\\\\\
vorgelegt am:  &  12. Mai 2011\\[1.8ex]

Studienbereich: &  \studienbereich\\
Fakult\"at: &  \mbox{Vermessung, Informatik und Mathematik}\\
Bildungseinrichtung: &  \mbox{Hochschule f\"ur Technik Stuttgart}\\[1.8ex]


von:            &  \autor\\[1.8ex]

Matrikelnummer: &  \matrikelnr\\[1.8ex]

Pr\"ufer:  &  \mbox{\erstgutachter, HfT Stuttgart}\\[1.8ex]


\end{tabular}


\begin{center}
\textbf{Zusammenfassung}

\begin{abstract}
In Zeiten in denen davon gesprochen wird, dass Moore's Law in absehbarer Zukunft
nicht mehr gilt werden zusehens Architekturen eingesetzt, die nicht mehr darauf
vertrauen, dass jede einzelne CPU ausreichend schnell ist, sondern vielmehr 
massiv parallele Rechenleistung bieten. Im Zuge dieses Architekturwandels müssen
gleichwertig die eingesetzten Software-Stacks überdacht und teilweise neu
ausgelegt werden, um die neue Form von Rechenleistung konsolodiert und effizient
zu nutzen.

Clojure ist ein LISP Dialekt und damit eine primär funktionale
Programmiersprache, die es sich zum Ziel genommen hat dem Entwickler und der
Anwendung hohe Parallelität und Performanz bei gleichzeitig hohem Komfort
bereitzustellen. Diese Arbeit bietet einen Überblick über die dahinter stehende
Methodik und grenzt sie gegen den bisherigen Verfahren ab
\end{abstract}


%\copyright\ \jahr\\[9ex]

\end{center}

\end{titlepage}

%\include{Inhalt/Abstract}
\ofoot{\pagemark}

% Vor dem Hauptteil: römische Ziffern
\pagenumbering{Roman}
\tableofcontents % Inhaltsverzeichnis

% Abkürzungsverzeichnis 
\nomenclature{XML}{Extensible Markup Language}

% für korrekte Überschrift in der Kopfzeile
\clearpage\markboth{\nomname}{\nomname} 
\printnomenclature
\label{sec:Glossar}

%\listoffigures % Abbildungsverzeichnis
\addcontentsline{toc}{chapter}{Abbildungsverzeichnis}
%\listoftables % Tabellenverzeichnis
\addcontentsline{toc}{chapter}{Tabellenverzeichnis}
\renewcommand{\lstlistlistingname}{Verzeichnis der Listings}
\addcontentsline{toc}{chapter}{Verzeichnis der Listings}
\lstlistoflistings % Listings-Verzeichnis

% arabische Seitenzahlen im Hauptteil 
\clearpage
\pagenumbering{arabic}

\chapter{Einleitung}
\label{cha:Einleitung}

\section{Motivation und Ziel der Arbeit}

Die vorliegende Arbeit ist die schriftliche Ausarbeitung des Seminar2 zum Thema
\NeuerBegriff{Clojure}. In ihr wird Clojure gegen�ber LISP abgegrenzt und �ber
die Vorteile dynamischer und funktionalier Programmierung im Gegenzug zu
konvetiell strukturierter Programmierung besprochen.

\section{Typographische Konventionen}

Folgende typographische Konventionen sind in dieser Arbeit eingesetzt.

\begin{itemize}
    \item 
        \NeuerBegriff{Neuer Begriff}\\
        Neue Begriffe sind f�r den schnellen �berblick gesondert im Textbild
        hervorgehoben.
    \item 
        \Fachbegriff{Fachbegriff}\\
        Fachbegriffe sind aus dem selben Grund wie neue Begriffe hervorgehoben.
    \item 
        \Eingabe{Eingabe}\\
        Referenzen auf Tastatureingaben sind als solche gekennzeichnet.
    \item 
        \Code{Quellcode}\\
        Quellcode kann wie eine Eingabe im Text eingebettet werden. Bei
        gr��eren Code-Versatzst�cken wird jedoch volles Syntax-Highlighting
        verwendet.
    \item 
        \Datei{C:\bs Pfad\bs Datei}\\
        Pfad- und Datei-Angaben
    \item 
        \Datentyp{Datentyp}\\
        Referenzen auf interne Datenstrukturen und Variablennamen tragen die
        typographische Kennzeichnung Datentyp.
\end{itemize}

%\include{Inhalt/Fazit}


\renewcommand{\bibname}{Quellenverzeichnis}
\bibliography{Bibliographie} % Aufruf: bibtex Masterarbeit
\addcontentsline{toc}{chapter}{\bibname}
%\bibliographystyle{alphadin} % DIN-Stil des Literaturverzeichnisses
\bibliographystyle{natdin} % DIN-Stil des Literaturverzeichnisses

% Anhang 
%\begin{appendix}
%    %\clearpage
%    \pagenumbering{roman}
%    \chapter{Anhang}
%    \label{sec:Anhang}
%    % Rand der Aufzählungen in Tabellen anpassen
%    \setdefaultleftmargin{1em}{}{}{}{}{}
%    \include{Inhalt/Anhang1}

%\end{appendix}

\newpage

% Index 
%\printindex

\end{document}
